\documentclass[../../main.tex]{subfiles}
\begin{document}

\begin{center}
  {\fontsize{48pt}{58pt}\selectfont \textbf{2007}\normalsize\textbf{年度}}
\end{center}

\begin{table}[htb]
  \centering
  \caption{出題テーマと難易度}
  \begin{tabular}{|c|p{9cm}|c|}
    \hline
    \textbf{問題番号} & \textbf{テーマ} & \textbf{難易度} \\
    \hline
    第1問 & 素数の冪による整除と組の個数 & \\
    \hline
    第2問 & 放物線の接線の回転と面積比の極限 & \\
    \hline
    第3問 & 正八角形上の点と三角形の面積の最大値 & \\
    \hline
    第4問 & 放物線と指数関数の接点，面積の和の極限 & \\
    \hline
  \end{tabular}
\end{table}

\begin{table}[htb]
  \centering
  \caption{解答}
  \renewcommand{\arraystretch}{1.4}
  \begin{tabular}{|c|c|p{9cm}|}
    \hline
    \textbf{問題番号} & \textbf{小問} & \textbf{答え} \\
    \hline
    第1問 & (1) & $p^{n+1-m}-p^{n-m}$個 \\
    \cline{2-3}
    & (2) & $(n+2)p^{n+1}-(n+1)p^n$組 \\
    \hline
    第2問 & (1) & $y=\dfrac{2\sqrt{3}a-1}{2a+\sqrt{3}}(x-a)+a^2$ \\
    \cline{2-3}
    & (2) & $\displaystyle\lim_{a\to\infty}\dfrac{T(a)}{S(a)}=4$ \\
    \hline
    第3問 & (1) & 最大値 $1+\dfrac{3\sqrt{2}}{4}$ \\
    \cline{2-3}
    & (2) & 最大値 $1+\dfrac{\sqrt{2}}{2}$ \\
    \hline
    第4問 & (1) & $a_n=(2+\sqrt{5})e^{\frac{1-\sqrt{5}}{2}}e^{-n}$ \\
    \cline{2-3}
    & (2) & $\displaystyle\lim_{n\to\infty}\sum_{k=0}^{n}S_k
    =1-\dfrac{(2+\sqrt{5})e^{\frac{3-\sqrt{5}}{2}}}{6(e-1)}$ \\
    \hline
  \end{tabular}
\end{table}

\paragraph{試験形態}
4問．試験時間は要確認．

\end{document}
